\documentclass[11pt,a4paper]{article}

% --- Typography & Geometry ---
\usepackage[utf8]{inputenc}
\usepackage[T1]{fontenc}
\usepackage{lmodern}
\usepackage{microtype}
\usepackage[margin=0.95in,headheight=15pt]{geometry}
\usepackage{setspace}
\setstretch{1.15}

% --- Math & Symbols ---
\usepackage{amsmath,amssymb,amsfonts,bm}

% --- Colors & Graphics ---
\usepackage{xcolor}
\usepackage{graphicx}
\usepackage{subcaption}
\usepackage{tikz}

% --- Earth & Copper / Botanical Color Palette ---
\definecolor{ForestGreen}{RGB}{28, 74, 43}      % Primary Headers & Concepts
\definecolor{Terracotta}{RGB}{184, 80, 58}      % Accent / Warnings / Contrasts
\definecolor{HarvestAmber}{RGB}{214, 130, 42}   % Tertiary Highlight
\definecolor{SageGreen}{RGB}{125, 162, 134}     % Subdued Rules & Borders
\definecolor{SlateBlue}{RGB}{43, 92, 143}       % Secondary Accent
\definecolor{LinenBg}{RGB}{252, 250, 246}       % Box Background
\definecolor{SoftSand}{RGB}{243, 239, 232}      % Header Tables
\definecolor{BodyText}{RGB}{40, 48, 44}         % Soft Charcoal Body

\usepackage[colorlinks=true,
            linkcolor=ForestGreen,
            citecolor=Terracotta,
            urlcolor=HarvestAmber]{hyperref}

% --- Tables ---
\usepackage{booktabs}
\usepackage{tabularx}
\usepackage{array}
\usepackage{colortbl}
\newcolumntype{Y}{>{\centering\arraybackslash}X}

% --- Boxes & Callouts ---
\usepackage[most]{tcolorbox}

\tcbset{
    sharp corners,
    boxrule=0.8pt,
    colback=LinenBg,
    colframe=SageGreen,
    fonttitle=\bfseries\color{ForestGreen},
    attach boxed title to top left={yshift=-2mm, xshift=4mm},
    boxed title style={colback=LinenBg, colframe=SageGreen, boxrule=0.6pt},
    top=12pt, bottom=10pt, left=12pt, right=12pt
}

\newtcolorbox{conceptbox}[1]{
    title={#1},
    colframe=ForestGreen,
    colbacktitle=LinenBg,
    fonttitle=\bfseries\color{ForestGreen}
}

\newtcolorbox{resultbox}[1]{
    title={#1},
    colframe=Terracotta,
    colbacktitle=LinenBg,
    fonttitle=\bfseries\color{Terracotta}
}

\newtcolorbox{highlightbox}[1]{
    title={#1},
    colframe=HarvestAmber,
    colbacktitle=LinenBg,
    fonttitle=\bfseries\color{HarvestAmber}
}

% --- Section Headings Styling ---
\usepackage{titlesec}
\titleformat{\section}
  {\Large\bfseries\color{ForestGreen}}{\thesection}{1em}{}[\titlerule]
\titleformat{\subsection}
  {\large\bfseries\color{ForestGreen}}{\thesubsection}{1em}{}
\titleformat{\subsubsection}
  {\normalsize\bfseries\color{Terracotta}}{\thesubsubsection}{1em}{}

% --- Headers & Footers ---
\usepackage{fancyhdr}
\pagestyle{fancy}
\fancyhf{}
\renewcommand{\headrulewidth}{0.4pt}
\renewcommand{\footrulewidth}{0.4pt}
\fancyhead[L]{\small\color{ForestGreen}\textbf{UAV Cotton Crop Pathology Classification}}
\fancyhead[R]{\small\color{BodyText}Multi-Architecture \& Transfer Learning Study}
\fancyfoot[C]{\color{BodyText}\thepage}

% --- Document Metadata ---
\title{\vspace{-1.5cm}\Huge\textbf{\color{ForestGreen}UAV-Assisted Cotton Crop Pathology Classification}\\[0.3em]
\Large\textbf{\color{Terracotta}A Comprehensive Empirical Investigation Across Convolutional Architectures,\\The Tripartite Transfer Learning Taxonomy, and Drone Flight Envelopes}}
\author{\textbf{Autonomous Agricultural Vision Engineering Laboratory}}
\date{\today}

\begin{document}
\maketitle

\begin{abstract}
\noindent
Precision agriculture leveraging Unmanned Aerial Vehicles (UAVs) demands computer vision architectures that reconcile high diagnostic precision with constrained embedded companion compute. This technical report presents a comprehensive empirical and theoretical investigation into convolutional neural network (CNN) architectures, inductive biases, and transfer learning mechanisms for classifying five cotton crop pathologies and abiotic stresses from drone aerial imagery. We first reveal and remediate severe data leakage embedded within the raw benchmark splits (including 1,040 cross-split duplicate images and class imbalance). We then evaluate a multi-model comparative suite encompassing a custom 4-stage Base CNN trained from scratch, MobileNetV3-Small (trained from scratch and fully fine-tuned), ResNet-18 evaluated across two distinct transfer learning paradigms (Linear Probing with a frozen backbone vs. Staged Partial Fine-Tuning), and compound-scaled EfficientNet-B3. Our findings demonstrate an extraordinary performance spectrum: while from-scratch baselines plateau at \textbf{71.51\%} (Base CNN) and \textbf{74.87\%} (MobileNetV3), freezing 99.98\% of ResNet-18 and training a 2,565-parameter linear probe immediately elevates test accuracy to \textbf{83.56\%} (81.98\% Macro F1). Partial fine-tuning of deep residual blocks further surges performance to \textbf{99.74\%} at \textbf{310.9 FPS}, while end-to-end transfer learning on MobileNetV3-Small achieves a top-tier \textbf{99.91\% test accuracy} and \textbf{99.88\% Macro F1} at \textbf{162.7 FPS} (6.15\,ms), matching the high-capacity EfficientNet-B3 (56.0 FPS) while consuming $7\times$ fewer parameters and delivering $2.9\times$ higher frame throughput.
\end{abstract}

\vspace{0.5cm}

% =========================================================================
\section{Introduction and Domain Challenges in UAV Vision}
% =========================================================================

Cotton (\textit{Gossypium hirsutum}) serves as an indispensable economic fiber globally, but is acutely vulnerable to an array of biotic pests and abiotic environmental disruptions. Early, non-destructive identification of canopy distress is vital to prevent catastrophic yield attrition. UAV-mounted multi-spectral and RGB aerial imaging represents a scalable modality for spatial crop surveillance; however, overhead drone photography introduces severe computer vision hurdles that fundamentally distinguish it from standard consumer image recognition:

\begin{enumerate}
    \item \textbf{Rotational Invariance \& Omnidirectional Perspectives}: Standard ground-level photographs feature a canonical upright orientation governed by the horizon. In contrast, aerial UAV camera sensors capture cotton canopies from nadir or steep oblique angles, making disease patterns completely orientation-agnostic.
    \item \textbf{Extreme Scale and Altitude Variations}: Canopy ground sampling distances (GSD) fluctuate continuously with drone flight altitude, wind shear, and terrain contours. Small boreholes drilled by bollworms appear across mere dozens of pixels, whereas generalized sunscald burns span hundreds of pixels.
    \item \textbf{Complex Co-Occurring Stressors}: Real-world open-field agriculture exhibits combined symptom complexes. For instance, plants suffering from root vascular wilt undergo compromised transpiration, making their leaves significantly more prone to secondary heat scorch and opportunistic insect infestation.
\end{enumerate}

\begin{figure}[htbp]
    \centering
    \includegraphics[width=\textwidth]{results/dataset_samples.png}
    \caption{\textbf{Visual Taxonomy of the Five Targeted Cotton Canopy Conditions}. From left to right: \textit{Cotton Bollworm Infestation} (localized boreholes and square feeding), \textit{Hail damage} (mechanical leaf lacerations and bruises), \textit{Mixed} (co-occurring biotic/abiotic trauma), \textit{Root Damage} (vascular wilting and yellowing chlorosis), and \textit{Strong-Light Burn} (sunscald necrosis and bronzing).}
    \label{fig:dataset_samples}
\end{figure}

% =========================================================================
\section{Dataset Forensics, Leakage Rectification, and Stratification}
% =========================================================================

\subsection{Forensic Audit of the Raw Dataset}
An exhaustive hash-level forensic analysis was performed on the raw image collection (\textit{Bangladesh-UAV-PlantDisease}), which contained 10,640 total files across three pre-arranged directories: \texttt{Training Dataset} (5,868 files), \texttt{Test Dataset} (4,772 files), and \texttt{Valid dataset} (747 files). 

Our automated audit uncovered critical data engineering flaws that would invalidate any naive machine learning evaluation:
\begin{itemize}
    \item \textbf{Catastrophic Data Leakage}: Exactly \textbf{1,040 identical images} (verified via cryptographic MD5 hashing) were present concurrently in both the raw training and test partitions. Evaluating models on the raw test set would yield deceptively inflated performance due to memorization of training instances.
    \item \textbf{Artificial Duplicate Inflation}: In \texttt{Test Dataset/Hail damages}, 509 unique physical images were duplicated across 1,988 files, distorting test distribution weights.
    \item \textbf{Unlabelled Validation Anomaly}: The raw \texttt{Valid dataset} directory contained 747 loose images without class subfolders; 647 of these were bit-for-bit duplicate copies of images in \texttt{Training Dataset/Mixed}.
    \item \textbf{Syntactic Inconsistency}: The class folder was designated \texttt{Hail damage} in training but \texttt{Hail damages} (plural) in testing.
\end{itemize}

\subsection{Sanitization Pipeline and Stratified Splitting}
To establish an uncompromised scientific benchmark, all 10,640 files were processed through an automated deduplication and normalization engine:
\begin{enumerate}
    \item Cryptographic MD5 content hashes were computed for every image file across all splits.
    \item Every unique hash was cross-referenced across class labels. Crucially, this revealed \textbf{zero label conflicts}---every unique image mapped deterministically to a single ground-truth pathology.
    \item The total clean, deduplicated universe yielded exactly \textbf{7,744 unique images}.
    \item A strict, deterministic \textbf{70\% / 15\% / 15\% Stratified Split} was generated using fixed seed initialization ($\text{seed}=42$), ensuring identical class proportions across training, validation, and test subsets.
\end{enumerate}

\begin{table}[htbp]
\centering
\small
\rowcolors{2}{LinenBg}{white}
\setlength{\tabcolsep}{5pt}
\begin{tabularx}{\textwidth}{c>{\raggedright\arraybackslash}Xcccc}
\toprule
\rowcolor{SoftSand}
\textbf{Class ID} & \textbf{Pathology / Stress Category} & \textbf{Total Clean Samples} & \textbf{Train (70\%)} & \textbf{Val (15\%)} & \textbf{Test (15\%)} \\
\midrule
0 & Cotton Bollworm Infestation & 1,033 (13.3\%) & 723 & 155 & 155 \\
1 & Hail damage & 2,655 (34.3\%) & 1,858 & 398 & 399 \\
2 & Mixed Distress & 2,394 (30.9\%) & 1,676 & 359 & 359 \\
3 & Root Damage & 908 (11.7\%) & 636 & 136 & 136 \\
4 & Strong-Light Burn & 754 (9.7\%) & 528 & 113 & 113 \\
\midrule
\rowcolor{SoftSand}
\textbf{Total} & \textbf{Sanitized Universe} & \textbf{7,744 (100.0\%)} & \textbf{5,421} & \textbf{1,161} & \textbf{1,162} \\
\bottomrule
\end{tabularx}
\caption{\textbf{Sanitized and Leakage-Free Dataset Stratification}. Every test instance is strictly disjoint by content hash from training and validation sets.}
\label{tab:splits}
\end{table}

% =========================================================================
\section{Convolutional Architectures and Computational Mechanics}
% =========================================================================

\subsection{Standard Convolution vs. Depthwise Separable Convolution}
Standard convolutional layers compute spatial filtering and channel cross-correlation concurrently, resulting in prohibitive computational footprints. Given an input feature tensor $X \in \mathbb{R}^{D_F \times D_F \times M}$ and an output tensor $Y \in \mathbb{R}^{D_F \times D_F \times N}$ with kernel dimension $D_K$, the computational complexity of standard 2D convolution is:
\begin{equation}
    \text{Cost}_{\text{Std}} = D_K \cdot D_K \cdot M \cdot N \cdot D_F \cdot D_F
\end{equation}

\textbf{Depthwise Separable Convolutions} decouple this process into two discrete operations:
\begin{enumerate}
    \item \textit{Depthwise Convolution}: Applies a single spatial filter per input channel:
    \begin{equation}
        \text{Cost}_{\text{Depthwise}} = D_K \cdot D_K \cdot M \cdot D_F \cdot D_F
    \end{equation}
    \item \textit{Pointwise ($1\times 1$) Convolution}: Computes a linear combination across all channels:
    \begin{equation}
        \text{Cost}_{\text{Pointwise}} = 1 \cdot 1 \cdot M \cdot N \cdot D_F \cdot D_F
    \end{equation}
\end{enumerate}

The resulting computational reduction ratio $\rho$ is mathematically defined as:
\begin{equation}
    \rho = \frac{\text{Cost}_{\text{Depthwise}} + \text{Cost}_{\text{Pointwise}}}{\text{Cost}_{\text{Std}}} = \frac{D_K^2 \cdot M \cdot D_F^2 + M \cdot N \cdot D_F^2}{D_K^2 \cdot M \cdot N \cdot D_F^2} = \frac{1}{N} + \frac{1}{D_K^2}
\end{equation}
For a typical $3 \times 3$ kernel ($D_K = 3$) with large channel depth $N$, $\rho \approx \frac{1}{9}$, cutting floating-point operations by nearly \textbf{89\%} with negligible representational degradation.

\subsection{Custom 4-Stage Base CNN: The Canonical From-Scratch Baseline}
To rigorously quantify the empirical lift attributable to transfer learning and advanced architectural motifs, we designed a canonical 4-stage convolutional baseline network (\texttt{CustomBaseCNN}) trained strictly from scratch without any external pretraining priors. Its architecture follows a classical hierarchical feature extraction topology:
\begin{itemize}
    \item \textbf{Stage 1 ($224\times 224 \to 112\times 112$)}: $\text{Conv2D}(3 \to 32, 3\times 3, \text{pad}=1) \to \text{BatchNorm2D} \to \text{ReLU} \to \text{MaxPool2D}(2, 2)$.
    \item \textbf{Stage 2 ($112\times 112 \to 56\times 56$)}: $\text{Conv2D}(32 \to 64, 3\times 3, \text{pad}=1) \to \text{BatchNorm2D} \to \text{ReLU} \to \text{MaxPool2D}(2, 2)$.
    \item \textbf{Stage 3 ($56\times 56 \to 28\times 28$)}: $\text{Conv2D}(64 \to 128, 3\times 3, \text{pad}=1) \to \text{BatchNorm2D} \to \text{ReLU} \to \text{MaxPool2D}(2, 2)$.
    \item \textbf{Stage 4 ($28\times 28 \to 14\times 14$)}: $\text{Conv2D}(128 \to 256, 3\times 3, \text{pad}=1) \to \text{BatchNorm2D} \to \text{ReLU} \to \text{MaxPool2D}(2, 2)$.
    \item \textbf{Classification Head}: $\text{AdaptiveAvgPool2D}(1, 1) \to \text{Flatten} \to \text{Linear}(256 \to 128) \to \text{ReLU} \to \text{Dropout}(p=0.4) \to \text{Linear}(128 \to 5)$.
\end{itemize}
All convolutional weights are initialized via Kaiming (He) normal initialization ($\mathcal{N}(0, \sqrt{2/n_{\text{in}}})$). Total parameter count is \textbf{422,437}, providing an ultra-compact baseline to establish the absolute lower bound of scratch learning capability.

\subsection{MobileNetV3-Small: The Ultra-Lightweight Challenger}
MobileNetV3-Small ($\approx 1.52$\,M parameters) synthesizes hardware-aware neural architecture search (NAS) with three pivotal innovations:

\begin{conceptbox}{Key Innovations in MobileNetV3-Small}
\begin{itemize}
    \item \textbf{Inverted Residuals with Linear Bottlenecks}: Features are first expanded to a high-dimensional space ($1\times 1$ conv), filtered through spatial depthwise convolutions ($3\times 3$ or $5\times 5$), and then projected back to a low-dimensional manifold without non-linearities, preserving manifold topology.
    \item \textbf{Squeeze-and-Excitation (SE) Dynamic Channel Attention}: SE blocks adaptively re-calibrate channel feature maps by capturing global inter-channel dependencies. Given feature tensor $U$, global average pooling produces spatial descriptor $z \in \mathbb{R}^C$:
    \begin{equation}
        s = \sigma\Big(W_2 \cdot \text{ReLU}\big(W_1 \cdot z\big)\Big)
    \end{equation}
    where $\sigma$ is the sigmoid function, and $W_1 \in \mathbb{R}^{\frac{C}{r} \times C}, W_2 \in \mathbb{R}^{C \times \frac{C}{r}}$ represent reduction and expansion projections. The channel outputs are then scaled: $\widetilde{X}_c = s_c \cdot U_c$.
    \item \textbf{Hardswish Non-Linearity}: Eliminates costly exponential computations associated with Swish ($\text{swish}(x) = x \cdot \sigma(x)$) on mobile hardware:
    \begin{equation}
        \text{h-swish}(x) = x \cdot \frac{\text{ReLU6}(x + 3)}{6} = x \cdot \frac{\min(\max(0, x + 3), 6)}{6}
    \end{equation}
\end{itemize}
\end{conceptbox}

\subsection{ResNet-18: Deep Residual Learning \& Identity Skip Mappings}
ResNet-18 ($\approx 11.18$\,M parameters) addresses the degradation and vanishing gradient problems in deep networks by introducing additive identity shortcut mappings. Formally, rather than forcing a stack of convolutional layers to directly approximate an underlying mapping $\mathcal{H}(x)$, the layers are parameterized to approximate a residual mapping:
\begin{equation}
    \mathcal{F}(x) = \mathcal{H}(x) - x \quad \implies \quad \mathcal{H}(x) = \mathcal{F}(x, \{W_i\}) + x
\end{equation}
During backpropagation, the gradient of the loss $\mathcal{E}$ with respect to the input activation $x$ is:
\begin{equation}
    \frac{\partial \mathcal{E}}{\partial x} = \frac{\partial \mathcal{E}}{\partial \mathcal{H}} \left( \frac{\partial \mathcal{F}}{\partial x} + 1 \right)
\end{equation}
The additive constant $+1$ ensures that gradients can propagate directly across residual blocks back to early layers without undergoing exponential attenuation, enabling rapid training convergence and stable representations.

\subsection{EfficientNet-B3: The Multi-Scale Compound Heavyweight}
EfficientNet-B3 ($\approx 10.70$\,M parameters) serves as the high-capacity benchmark. Rather than arbitrarily tuning depth, width, or resolution in isolation, EfficientNet relies on the \textbf{Compound Scaling Principle}.

\begin{highlightbox}{The Compound Scaling Theorem}
Given a target computational resource expansion factor $\alpha \cdot \beta^2 \cdot \gamma^2 \approx 2$, network depth $d$, channel width $w$, and input image resolution $r$ are scaled concurrently using compound coefficient $\phi$:
\begin{equation}
    d = \alpha^\phi, \quad w = \beta^\phi, \quad r = \gamma^\phi
\end{equation}
For EfficientNet-B3 ($\phi = 3$):
\begin{itemize}
    \item \textbf{Input Resolution}: $r = 300 \times 300$ pixels (retaining high Nyquist spatial frequency).
    \item \textbf{Backbone}: 26 stacked MBConv blocks utilizing Swish activations, Squeeze-and-Excitation modules with squeeze ratio 0.25, and Stochastic Depth (drop-connect rate 0.3) for strong regularization.
\end{itemize}
\end{highlightbox}

% =========================================================================
\section{Transfer Learning Theory and the Tripartite Taxonomy}
% =========================================================================

\subsection{Mathematical Formalization of Transfer Learning}
Formally, a domain $\mathcal{D} = \{\mathcal{X}, P(X)\}$ consists of a feature space $\mathcal{X}$ and a marginal probability distribution $P(X)$, where $X = \{x_1, \dots, x_n\} \subset \mathcal{X}$. A task $\mathcal{T} = \{\mathcal{Y}, f(\cdot)\}$ consists of a label space $\mathcal{Y}$ and an objective predictive function $f(\cdot)$ that maps inputs to posterior probabilities $P(y|x)$.

In our study:
\begin{itemize}
    \item \textbf{Source Domain \& Task ($\mathcal{D}_S, \mathcal{T}_S$)}: $\mathcal{X}_S = \mathbb{R}^{224 \times 224 \times 3}$ encompassing 1.28 million natural images from the ImageNet-1K challenge with 1,000 diverse classes ($\mathcal{Y}_S = \{1, \dots, 1000\}$).
    \item \textbf{Target Domain \& Task ($\mathcal{D}_T, \mathcal{T}_T$)}: $\mathcal{X}_T$ consists of overhead UAV aerial cotton crops, with label space $\mathcal{Y}_T = \{0, 1, 2, 3, 4\}$ representing the 5 specific pathology states.
\end{itemize}

Transfer learning exploits the condition that the early conditional distributions of low-level visual features are highly invariant across natural domains:
\begin{equation}
    P_{\mathcal{D}_S}(\text{edges}, \text{textures} \mid X_S) \approx P_{\mathcal{D}_T}(\text{edges}, \text{textures} \mid X_T)
\end{equation}

\subsection{The Tripartite Taxonomy of Transfer Learning Paradigms}
To systematically investigate how transferred knowledge interacts with aerial pathology features, we formulate and test three distinct transfer learning paradigms:

\begin{enumerate}
    \item \textbf{Paradigm 1: Fixed Feature Extraction / Linear Probing (Linear Probe)}:
    The entire convolutional feature backbone $\Phi(\cdot; \Theta_{\text{frozen}})$ is strictly locked ($\nabla_{\Theta} \mathcal{L} = 0$). Only the newly initialized linear classification projection matrix $W \in \mathbb{R}^{C \times d}$ and bias $b \in \mathbb{R}^C$ are optimized:
    \begin{equation}
        \min_{W, b} \frac{1}{|\mathcal{B}|} \sum_{i \in \mathcal{B}} \mathcal{L}_{\text{CW}}\Big(W \cdot \Phi(x_i; \Theta_{\text{frozen}}) + b, \, y_i\Big)
    \end{equation}
    In our experiment on \textbf{ResNet-18}, exactly \textbf{2,565 parameters} (0.02\% of the network) are trained, while 11,176,512 parameters remain locked. This rigorously tests whether generic ImageNet representations are linearly separable for crop diseases without any domain adaptation.
    
    \item \textbf{Paradigm 2: Staged / Partial Layer-Wise Fine-Tuning (Partial Transfer)}:
    The early visual stages $\Theta_{\text{shallow}} = \{\text{conv1}, \text{bn1}, \text{layer1}, \text{layer2}\}$ are frozen to lock generic edge and texture primitives, while the deep semantic stages $\Theta_{\text{deep}} = \{\text{layer3}, \text{layer4}, \text{fc}\}$ are unlocked for gradient updates:
    \begin{equation}
        \Theta_{\text{shallow}}^{(t+1)} = \Theta_{\text{shallow}}^{(t)}, \qquad \Theta_{\text{deep}}^{(t+1)} = \Theta_{\text{deep}}^{(t)} - \eta_t \nabla_{\Theta_{\text{deep}}} \mathcal{L}_{\text{CW}}
    \end{equation}
    In our \textbf{ResNet-18 Partial} implementation, \textbf{6,069,253 parameters (54.3\%)} are actively trained. This allows higher-order receptive fields to specialize for necrotic tissue, wilting, and insect boreholes while preserving foundational spatial filters.
    
    \item \textbf{Paradigm 3: Full End-to-End Fine-Tuning (Full Transfer)}:
    All parameters across the entire network $\Theta = \{\Theta_{\text{backbone}}, \Theta_{\text{head}}\}$ are jointly updated under a conservative learning rate with cosine decay:
    \begin{equation}
        \Theta^{(t+1)} = \Theta^{(t)} - \eta_t \nabla_{\Theta} \mathcal{L}_{\text{CW}}
    \end{equation}
    Evaluated on \textbf{MobileNetV3-Small} (1.52M parameters) and \textbf{EfficientNet-B3} (10.70M parameters).
\end{enumerate}

% =========================================================================
\section{Loss Functions, Imbalance Balancing, and Training Strategy}
% =========================================================================

\subsection{Balanced Class-Weighted Cross-Entropy Loss}
As detailed in Table~\ref{tab:splits}, our cotton dataset exhibits substantial class imbalance: \textit{Hail damage} (34.3\%) and \textit{Mixed} (30.9\%) represent nearly two-thirds of the data, while \textit{Strong-Light Burn} comprises only 9.7\%. Standard cross-entropy would bias the model toward the majority classes.

To guarantee equitable gradient flow across minority conditions, we formulated a dynamically balanced class weighting vector $w \in \mathbb{R}^C$:
\begin{equation}
    w_c = \frac{N}{C \cdot N_c}, \quad \text{normalized such that} \quad \frac{1}{C}\sum_{c=1}^C w_c = 1.0
\end{equation}
where $N = 5,421$ is the total training population, $C=5$ is the number of classes, and $N_c$ is the training sample count for class $c$. The calculated weight vector is:
\begin{equation}
    w = \begin{bmatrix} 1.16 & 0.45 & 0.50 & 1.31 & 1.58 \end{bmatrix}^T
\end{equation}
The weighted multi-class cross-entropy loss for mini-batch $\mathcal{B}$ is computed as:
\begin{equation}
    \mathcal{L}_{\text{CW}} = -\frac{1}{|\mathcal{B}|} \sum_{i \in \mathcal{B}} w_{y_i} \log \left( \frac{\exp(z_{i, y_i})}{\sum_{j=1}^C \exp(z_{i, j})} \right)
\end{equation}
where $z_i$ represents the raw unnormalized logits output by the classifier head.

\subsection{Optimization, Mixed Precision, and Scheduling}
\begin{itemize}
    \item \textbf{Decoupled Weight Decay Optimizer (AdamW)}: Employed with initial base learning rate $\eta_0 = 3 \times 10^{-4}$ (or $\eta_0 = 10^{-3}$ for the linear probe) and weight decay coefficient $\lambda = 10^{-4}$.
    \item \textbf{Cosine Annealing Schedule}: Decays the learning rate to prevent late-stage gradient chatter:
    \begin{equation}
        \eta_t = \eta_{\min} + \frac{1}{2}(\eta_0 - \eta_{\min})\left(1 + \cos\left(\frac{t}{T_{\max}}\pi\right)\right), \quad \text{with } \eta_{\min} = 10^{-6}, \, T_{\max} = 10
    \end{equation}
    \item \textbf{Hardware Acceleration via PyTorch AMP}: Automated Mixed Precision (AMP) was utilized on the NVIDIA GTX 1650 (Turing architecture), reducing memory traffic by $\approx 50\%$ while maintaining numerical stability.
    \item \textbf{Drone-Perspective Data Augmentation}: Implemented $360^\circ$ random rotations, horizontal/vertical flips, random resized cropping (scale 0.75--1.0), and color jittering ($\pm 20\%$ brightness/contrast/saturation) to simulate varying sun angles and drone flight headings.
\end{itemize}

% =========================================================================
\section{Experimental Results and Empirical Comparative Analysis}
% =========================================================================

\subsection{Overall Model Scorecard}
All six experimental configurations were evaluated on the 1,162 strictly held-out test images. Table~\ref{tab:scorecard} summarizes the quantitative findings across model capacity, transfer paradigms, diagnostic precision, and inference throughput.

\begin{table}[htbp]
\centering
\rowcolors{2}{LinenBg}{white}
\small
\resizebox{\textwidth}{!}{%
\begin{tabular}{lcccccc}
\toprule
\rowcolor{SoftSand}
\textbf{Model \& Paradigm} & \textbf{Initial Weights} & \textbf{Total Params} & \textbf{Trainable} & \textbf{Test Acc.} & \textbf{Macro F1} & \textbf{Throughput (Latency)} \\
\midrule
Base CNN (Scratch) & Random (He) & 0.42 M & 422 K (100\%) & 71.51\% & 68.83\% & \textbf{826.4 FPS} (1.21 ms) \\
MobileNetV3-Small (Scratch) & Random (He) & 1.52 M & 1.52 M (100\%) & 74.87\% & 73.76\% & 164.4 FPS (6.08 ms) \\
\midrule
ResNet-18 (Linear Probe) & ImageNet-1K & 11.18 M & \textbf{2.5 K (0.02\%)} & 83.56\% & 81.98\% & 316.2 FPS (3.16 ms) \\
ResNet-18 (Partial Transfer) & ImageNet-1K & 11.18 M & 6.07 M (54.3\%) & 99.74\% & 99.68\% & 310.9 FPS (3.22 ms) \\
\textbf{MobileNetV3 (Full Transfer)} & \textbf{ImageNet-1K} & \textbf{1.52 M} & \textbf{1.52 M (100\%)} & \textbf{99.91\%} & \textbf{99.88\%} & \textbf{162.7 FPS (6.15 ms)} \\
EfficientNet-B3 (Scaled Transfer) & ImageNet-1K & 10.70 M & 10.70 M (100\%) & \textbf{99.91\%} & \textbf{99.88\%} & 56.0 FPS (17.86 ms) \\
\bottomrule
\end{tabular}%
}
\caption{\textbf{Master Performance Benchmark on 1,162 Held-Out UAV Test Images}. The study demonstrates the progression from scratch learning (71--75\%) to linear probing (83.56\%), partial transfer (99.74\%), and full transfer (99.91\%).}
\label{tab:scorecard}
\end{table}

\begin{figure}[htbp]
    \centering
    \includegraphics[width=0.98\textwidth]{results/models_comparison.png}
    \caption{\textbf{Quantitative Benchmark Comparison Across All Models}. Left: Test Accuracy and Macro F1 comparison highlighting the progressive gains from scratch baselines through Linear Probing, Partial Fine-Tuning, and Full Transfer Learning. Right: Edge inference latency on NVIDIA GTX 1650 (lower is better).}
    \label{fig:comparison}
\end{figure}

\subsection{The Progressive Spectrum of Transfer Learning}
Analyzing the empirical trajectory across paradigms reveals fundamental insights into deep convolutional representations:

\begin{resultbox}{Key Empirical Findings from the Transfer Spectrum}
\begin{itemize}
    \item \textbf{The Linear Probing Baseline (+12.05\% over Scratch)}: By freezing 99.98\% of ResNet-18 and training only a single linear layer (2,565 parameters), test accuracy reaches \textbf{83.56\%}---vastly exceeding the from-scratch Base CNN (71.51\%) and MobileNetV3 (74.87\%). This proves that ImageNet visual manifolds provide rich generic visual primitives (chlorotic edges, textural gradients) that are linearly separable even before domain fine-tuning.
    \item \textbf{The Partial Fine-Tuning Leap (+16.18\% over Linear Probe)}: Unfreezing the deep residual stages ($\text{layer3}, \text{layer4}$) allows the network to adapt semantic receptive fields, driving accuracy from 83.56\% to \textbf{99.74\%} and Macro F1 to \textbf{99.68\%}, while delivering over \textbf{310 FPS}.
    \item \textbf{Full Transfer Parity}: MobileNetV3-Small with full transfer achieves \textbf{99.91\% accuracy} (1,161 out of 1,162 correct), exactly matching the heavy EfficientNet-B3 while running at \textbf{162.7 FPS} versus EfficientNet's 56.0 FPS ($2.9\times$ speedup).
\end{itemize}
\end{resultbox}

\begin{figure}[htbp]
    \centering
    \includegraphics[width=\textwidth]{results/training_curves.png}
    \caption{\textbf{Validation Loss and Accuracy Trajectories Across 10 Epochs}. Pretrained models (solid lines) converge rapidly by Epoch 2--3, whereas from-scratch models (dashed lines) lag significantly throughout the optimization envelope.}
    \label{fig:curves}
\end{figure}

\subsection{Per-Class Metrics and Diagnostic Dissection}
Table~\ref{tab:per_class} compares per-class diagnostic performance across models on all 1,162 test samples.

\begin{table}[htbp]
\centering
\scriptsize
\rowcolors{2}{LinenBg}{white}
\resizebox{\textwidth}{!}{%
\begin{tabular}{l|ccc|ccc|ccc|c}
\toprule
\rowcolor{SoftSand}
& \multicolumn{3}{c|}{\textbf{Base CNN (Scratch)}} & \multicolumn{3}{c|}{\textbf{ResNet-18 (Linear Probe)}} & \multicolumn{3}{c|}{\textbf{MobileNetV3 (Transfer)}} & \\
\rowcolor{SoftSand}
\textbf{Pathology Category} & \textbf{Prec.} & \textbf{Rec.} & \textbf{F1} & \textbf{Prec.} & \textbf{Rec.} & \textbf{F1} & \textbf{Prec.} & \textbf{Rec.} & \textbf{F1} & \textbf{Support} \\
\midrule
Cotton Bollworm Infestation & 61.46\% & 81.29\% & 70.00\% & 93.33\% & 81.29\% & 86.90\% & 100.00\% & 100.00\% & \textbf{100.00\%} & 155 \\
Hail damage & 80.10\% & 76.69\% & 78.36\% & 91.48\% & 80.70\% & 85.75\% & 100.00\% & 100.00\% & \textbf{100.00\%} & 399 \\
Mixed Distress & 79.53\% & 66.02\% & 72.15\% & 81.77\% & 89.97\% & 85.68\% & 100.00\% & 99.72\% & \textbf{99.86\%} & 359 \\
Root Damage & 73.73\% & 63.97\% & 68.50\% & 70.83\% & 75.00\% & 72.86\% & 100.00\% & 100.00\% & \textbf{100.00\%} & 136 \\
Strong-Light Burn & 47.17\% & 66.37\% & 55.15\% & 72.06\% & 86.73\% & 78.71\% & 99.12\% & 100.00\% & \textbf{99.56\%} & 113 \\
\midrule
\rowcolor{SoftSand}
\textbf{Macro Average} & \textbf{68.40\%} & \textbf{70.87\%} & \textbf{68.83\%} & \textbf{81.89\%} & \textbf{82.74\%} & \textbf{81.98\%} & \textbf{99.82\%} & \textbf{99.94\%} & \textbf{99.88\%} & \textbf{1,162} \\
\bottomrule
\end{tabular}%
}
\caption{\textbf{Per-Class Performance Dissection on the Held-Out Test Set}. Notice the progression in minority class precision for Strong-Light Burn: from 47.17\% (Scratch) $\to$ 72.06\% (Linear Probe) $\to$ 99.12\% (Full Transfer).}
\label{tab:per_class}
\end{table}

\begin{figure}[htbp]
    \centering
    \begin{subfigure}[b]{0.32\textwidth}
        \includegraphics[width=\textwidth]{results/custom_base_cnn_scratch_confusion_matrix.png}
        \caption{Base CNN (Scratch: 71.51\%)}
    \end{subfigure}
    \hfill
    \begin{subfigure}[b]{0.32\textwidth}
        \includegraphics[width=\textwidth]{results/mobilenet_v3_small_scratch_confusion_matrix.png}
        \caption{MobileNetV3 (Scratch: 74.87\%)}
    \end{subfigure}
    \hfill
    \begin{subfigure}[b]{0.32\textwidth}
        \includegraphics[width=\textwidth]{results/resnet18_frozen_confusion_matrix.png}
        \caption{ResNet-18 (Linear Probe: 83.56\%)}
    \end{subfigure}
    \vskip\baselineskip
    \begin{subfigure}[b]{0.32\textwidth}
        \includegraphics[width=\textwidth]{results/resnet18_partial_confusion_matrix.png}
        \caption{ResNet-18 (Partial: 99.74\%)}
    \end{subfigure}
    \hfill
    \begin{subfigure}[b]{0.32\textwidth}
        \includegraphics[width=\textwidth]{results/mobilenet_v3_small_transfer_confusion_matrix.png}
        \caption{MobileNetV3 (Transfer: 99.91\%)}
    \end{subfigure}
    \hfill
    \begin{subfigure}[b]{0.32\textwidth}
        \includegraphics[width=\textwidth]{results/efficientnet_b3_transfer_confusion_matrix.png}
        \caption{EfficientNet-B3 (Transfer: 99.91\%)}
    \end{subfigure}
    \caption{\textbf{Comprehensive Confusion Matrix Suite on 1,162 Held-Out UAV Test Samples}. From scratch baselines (top left/center) display prominent off-diagonal misclassifications; linear probing (top right) consolidates predictions along the diagonal; partial and full transfer models (bottom row) achieve near-perfect classification (up to 1,161/1,162 correct).}
    \label{fig:confusion_matrices}
\end{figure}

% =========================================================================
\section{Edge Deployment and Drone Engineering Discussion}
% =========================================================================

\subsection{Real-Time Flight Envelope and Companion Compute}
In commercial UAV agricultural missions, cameras acquire imagery at 30 to 60 FPS. An effective onboard diagnostic payload must classify incoming frames in real-time without buffering backlogs:
\begin{itemize}
    \item \textbf{Base CNN Scratch}: While operating at \textbf{826.4 FPS} (1.21 ms), its 71.51\% diagnostic accuracy is unacceptable for autonomous spot spraying where misclassifying healthy crops or missing pests leads to crop damage.
    \item \textbf{ResNet-18 Partial Transfer}: Delivers \textbf{310.9 FPS} (3.22 ms) at \textbf{99.74\% accuracy}. For industrial companion computers with hardware Tensor Cores, ResNet-18 is exceptionally fast and simple to deploy via TensorRT.
    \item \textbf{MobileNetV3-Small Full Transfer}: The definitive \textbf{Pareto-optimal champion}. With \textbf{99.91\% accuracy}, \textbf{162.7 FPS} (6.15 ms/frame), and only \textbf{1.52M parameters}, it occupies only $37\%$ of a standard 60\,FPS frame budget (16.6\,ms). This leaves ample headroom for concurrent visual odometry, SLAM, and GPS geotagging on low-power platforms like the NVIDIA Jetson Orin Nano (7W--15W) or Raspberry Pi 5.
    \item \textbf{EfficientNet-B3}: At 17.86 ms on desktop GPU, throughput is limited to \textbf{56.0 FPS}, which risks dropping frames on edge devices without offering any diagnostic advantage over MobileNetV3.
\end{itemize}

% =========================================================================
\section{Conclusion}
% =========================================================================

This investigation resolved the critical data leakage vulnerabilities in the UAV cotton pathology benchmark and presented a comprehensive empirical evaluation of convolutional architectures and transfer learning paradigms. Specifically:
\begin{enumerate}
    \item \textbf{Leakage Elimination}: Purged 1,040 cross-split duplicates, creating a clean 7,744-image benchmark with zero label ambiguity.
    \item \textbf{The Transfer Learning Spectrum}: Established the quantitative hierarchy: Base CNN Scratch (71.51\%) $\to$ MobileNetV3 Scratch (74.87\%) $\to$ ResNet-18 Linear Probe (83.56\%) $\to$ ResNet-18 Partial Transfer (99.74\%) $\to$ MobileNetV3 Full Transfer (99.91\%).
    \item \textbf{Efficiency Verification}: Proved that low-parameter models with rich inductive transfer (MobileNetV3-Small) achieve full diagnostic parity with architectures $7\times$ their size (EfficientNet-B3) while operating at \textbf{162.7 FPS}.
\end{enumerate}

All model checkpoints, training histories, split manifests, and evaluation artifacts are fully preserved in the \texttt{TrenUAV} repository for edge compilation and onboard UAV deployment.

\end{document}
